\documentclass[12pt,a4paper]{extarticle}
\usepackage[utf8]{inputenc}
\usepackage[T5]{fontenc}
\usepackage{geometry}
\geometry{a4paper, left=1.5cm, right=1cm, top=1.5cm, bottom=1.5cm, headheight=20pt, headsep=15pt, footskip=28pt}
\usepackage{enumitem}
\usepackage{amsmath}
\usepackage{amssymb}
\usepackage{diagbox}
\usepackage[most]{tcolorbox}
\newcommand{\khongxacdinh}{\text{KXĐ}}
\newcommand{\blackbox}[1]{\texorpdfstring{\fcolorbox{black}{black}{\textcolor{white}{#1}}}{#1}}
\usepackage{tikz}
\definecolor{bookmagenta}{RGB}{58,90,172}
\definecolor{book}{RGB}{58,90,172}
\definecolor{myred}{RGB}{200,40,40}
\newcommand{\bluecheck}{\textcolor{myblue}{\checkmark}}
\newtcolorbox{formulaBox}{colback=gray!5,colframe=bookmagenta,boxrule=0.8pt,arc=2mm}
\newtcolorbox{notebox}{colback=blue!3,colframe=myblue,boxrule=0.8pt,arc=2mm}
\usetikzlibrary{patterns, patterns.meta} % Thêm patterns.meta vào đây
% Gọi package nằm trong thư mục con template
\usepackage{../template/mngocbox}

\begin{document}

% Sử dụng ngoặc vuông [...] để thay đổi linh hoạt các thông số
\makeheadbox[headerright={Tổng ôn 10-11},chude={CHỦ ĐỀ 10},tieude={Hình học không gian},dinhdang={Hình học},lop={12 },gv={Nguyễn Lê Minh Ngọc}]

\vspace{2em}
\makeheading{II}{Bài tập tự luận}
\textbf{Câu 1:} Cho hình chóp S.ABCD, đáy ABCD có AB cắt CD tại E, AC cắt BD tại F.
\begin{enumerate}
\item Tìm giao tuyến giữa các mặt phẳng (SAB) và (SCD), (SAC) và (SBD)
\item Tìm giao tuyến của (SEF) với các mặt phẳng (SAD), (SBC)
\end{enumerate}
\textbf{Câu 2:} Cho hình chóp S.ABCD. M là điểm nằm trên cạnh SC.
\begin{enumerate}
\item Tìm giao điểm của AM và (SBD)
\item Gọi N là một điểm nằm trên cạnh BC. Tìm giao điểm của SD và (AMN)
\end{enumerate}
\textbf{Câu 3: }Cho tứ diện ABCD. Gọi I, J lần lượt là trung điểm của các cạnh BC, CD. Trên cạnh AC lấy điểm K. Gọi M là giao điểm BK và AI, N là giao điểm của DK và AJ. Chứng minh rằng đường thẳng MN song song với đường thẳng BD\\
\vspace{0.5cm}
\textbf{Câu 4: }Cho tứ diện ABCD. Gọi M, N lần lượt là trọng tâm của tam giác ABC và BCD. Chứng minh rằng MN//(ABD) và MN//(ACD) \\ 
\vspace{0.5cm}
\textbf{Câu 5: }Cho hình hộp chữ nhật ABCD.A'B'C'D' có các cạnh bên song song với nhau.
\begin{enumerate}
\item Chứng minh rằng mặt phẳng (OMN) và (SBC) song song với nhau
\item Giả sử hai tam giác SAD và ABC đều là tam giác cân tại A. Gọi AE và AF lần lượt là các đường phân giác trong của tam giác ACD và SAB. Chứng minh EF song song với (SAD)
\end{enumerate}
\textbf{Câu 6: }Cho hình chóp S.ABCD có ABCD là hình vuông tâm O và SA vuông góc (ABCD). Gọi H, K lần lượt là hình chiếu vuông góc của điểm A trên các cạnh SB, SD.
\begin{enumerate}
\item BC vuông góc (SAB)
\item SC vuông góc (AHK)
\end{enumerate}
\textbf{Câu 7: } Cho hình chóp S.ABCD có đáy là hình thoi cạnh a, SA = $\sqrt(3)$ và SA vuông BC. Góc giữa SD và BC = ? \\ 
\vspace{0.5cm}
\textbf{Câu 8: }Cho hình chóp S.ABC. SA=2AB=2a và SA vuông góc với (ABC) khoảng cách từ A đến SB = ? \\ 
\vspace{0.5cm}
\textbf{Câu 9: }Cho hình chóp đều S.ABC và M,N,E lần lượt là trung điểm của các cạnh SA, SB, SC. Tính thể tích khối chóp cụt MNE.ABC. Biết thể tích S.ABC=60 $60m^3$\\





\end{document}
